
\section{Conditioning the interpolant path: proofs}
\label{appendix:proof_conditioning}

This appendix states and proves the posterior-path result used in Section~\ref{subsec:posterior_dynamics} and proves Theorem~\ref{theorem:unified_posterior}, the two facts that turn the prior interpolant path into the unified posterior sampler. Both reuse the structural results of Appendix~\ref{appendix:score_velocity}: the velocity--score duality \eqref{eq:general_velocity_of_score} and the marginal-preserving SDE family (Proposition~\ref{prop:marginal_preservation}).

\subsection{Posterior path}

\begin{proposition}[Posterior path]\label{prop:conditional_path}
Conditioning the interpolant \eqref{eq:general_interpolant} on $\obs$ yields the path
\begin{align}\label{eq:conditional_path}
    \bx_\tau = \alpha_\tau\bx_0 + \sigpath_\tau\latentnoise + \beta_\tau\bx_1,
    \qquad \bx_1\sim\conddist{\bx_1}{\obs,\bx_0},\ \ \latentnoise\sim\mathcal{N}(0,\bI),
\end{align}
with the same schedules $(\alpha_\tau,\beta_\tau,\sigpath_\tau)$ and anchor $\bx_0$, $\latentnoise$ still standard normal and independent of $\bx_1$, and target retargeted to the posterior. Its score and velocity are the prior ones shifted by the likelihood score,
\begin{align}
    \genscore_\tau^{\obs}(\bx,\bx_0)
    &= \genscore_\tau(\bx,\bx_0) + \nabla_{\bx}\log\Phi_\tau^{\obs}(\bx,\bx_0),
       \label{eq:conditional_score}\\
    \fmvel_\tau^{\obs}(\bx,\bx_0)
    &= \fmvel_\tau(\bx,\bx_0) + \vscoef_\tau\,\nabla_{\bx}\log\Phi_\tau^{\obs}(\bx,\bx_0),
       \label{eq:conditional_velocity}
\end{align}
with $\vscoef_\tau$ the velocity--score coefficient of Eq.~\eqref{eq:general_score_velocity} and $\Phi_\tau^{\obs}$ the likelihood tilt \eqref{eq:likelihood_tilt}.
\end{proposition}

\begin{proof}
Take the prior interpolant \eqref{eq:general_interpolant}, $\bx_\tau = \alpha_\tau\bx_0 + \sigpath_\tau\latentnoise + \beta_\tau\bx_1$, with $\latentnoise\sim\mathcal{N}(0,\bI)$ independent of $\bx_1$, and the observation $\obs = \obsoperator(\bx_1) + \obsnoise$ with $\obsnoise$ independent of $(\bx_1,\latentnoise)$. Conditioning the joint law of $(\bx_1,\latentnoise)$ on $\obs$ leaves the Gaussian driver unchanged: since $\latentnoise$ is independent of $(\bx_1,\obsnoise)$ and $\obs$ is a function of $(\bx_1,\obsnoise)$, we have $\latentnoise \perp \obs \mid \bx_1$, hence
\begin{align}
    \conddist{\latentnoise}{\bx_1,\obs} = \conddist{\latentnoise}{\bx_1} = \mathcal{N}(0,\bI).
\end{align}
So $\latentnoise\mid\obs\sim\mathcal{N}(0,\bI)$ and remains independent of $\bx_1\mid\obs$, giving \eqref{eq:conditional_path}: the same schedules and anchor $\bx_0$, with the target retargeted to $\conddist{\bx_1}{\obs,\bx_0}$. The path is therefore conditionally Gaussian given $\bx_1$, as in the prior.

For the score, Bayes' rule on the marginals gives
\begin{align}\label{eq:posterior_marginal}
    p_\tau^{\obs}(\bx\mid\bx_0) = \frac{\Phi_\tau^{\obs}(\bx,\bx_0)\,p_\tau(\bx\mid\bx_0)}{\conddist{\obs}{\bx_0}},
    \qquad
    \Phi_\tau^{\obs}(\bx,\bx_0)=\E[\conddist{\obs}{\bx_1}\mid\bx_\tau=\bx,\bx_0],
\end{align}
where the second equality uses $\obs\perp\bx_\tau\mid\bx_1$. Taking $\nabla_{\bx}\log$ yields \eqref{eq:conditional_score}. For the velocity, apply the duality \eqref{eq:general_velocity_of_score} to both paths. The posterior path shares the schedules and anchor of the prior, so the affine-in-$\bx$ part of \eqref{eq:general_velocity_of_score} is identical,
\begin{align}
    \fmvel_\tau^{\obs}(\bx)
    &= \tfrac{\dot\beta_\tau}{\beta_\tau}\bx + \big(\dot\alpha_\tau - \tfrac{\dot\beta_\tau\alpha_\tau}{\beta_\tau}\big)\bx_0 + \vscoef_\tau\,\genscore_\tau^{\obs}(\bx), \\
    \fmvel_\tau(\bx)
    &= \tfrac{\dot\beta_\tau}{\beta_\tau}\bx + \big(\dot\alpha_\tau - \tfrac{\dot\beta_\tau\alpha_\tau}{\beta_\tau}\big)\bx_0 + \vscoef_\tau\,\genscore_\tau(\bx),
\end{align}
and subtracting with \eqref{eq:conditional_score} gives \eqref{eq:conditional_velocity},
\begin{align}
    \fmvel_\tau^{\obs}(\bx) - \fmvel_\tau(\bx)
    = \vscoef_\tau\big(\genscore_\tau^{\obs}(\bx) - \genscore_\tau(\bx)\big)
    = \vscoef_\tau\,\nabla_{\bx}\log\Phi_\tau^{\obs}(\bx,\bx_0). \qedhere
\end{align}
\end{proof}

\subsection{Proof of Theorem~\ref{theorem:unified_posterior}}
\begin{proof}
By Proposition~\ref{prop:conditional_path} the posterior marginals $p_\tau^{\obs}$ are those of an interpolant path with velocity $\fmvel_\tau^{\obs}$ and score $\genscore_\tau^{\obs}$. Applying the marginal-preserving SDE family of Proposition~\ref{prop:marginal_preservation} to this path, for any $\gdiff_\tau\ge0$ the SDE
\begin{align}
    \rd\bXstar_\tau = \left[\fmvel_\tau^{\obs}(\bXstar_\tau,\bx_0) + \tfrac12\gdiff_\tau^2\,\genscore_\tau^{\obs}(\bXstar_\tau,\bx_0)\right]\rd\tau + \gdiff_\tau\,\rd\bW_\tau,
\end{align}
started from $\bXstar_0\sim p_0^{\obs}$, has marginal law $p_\tau^{\obs}$ for all $\tau\in[0,1]$; in particular $\bXstar_1\sim p_1^{\obs}=\conddist{\bx_1}{\obs,\bx_0}$. By \eqref{eq:posterior_marginal} the initial law is $p_0^{\obs}(\rd\bx)\propto\Phi_0^{\obs}(\bx,\bx_0)\,p_0(\rd\bx)$. Substituting \eqref{eq:conditional_score}--\eqref{eq:conditional_velocity} gives the drift-correction form \eqref{eq:unified_posterior_drift} with guidance weight $\gweight_\tau=\vscoef_\tau+\tfrac12\gdiff_\tau^2$.

The initial reweighting is vacuous for both model classes. Because $\beta_0=0$, the source $\bx_{\tau=0}=\alpha_0\bx_0+\sigpath_0\latentnoise$ does not depend on $\bx_1$, so under the coupling $\latentnoise\perp\bx_1$ the tilt $\Phi_0^{\obs}(\bx,\bx_0)=\conddist{\obs}{\bx_0}$ is constant in $\bx$. Hence $p_0^{\obs}=p_0$: for flow matching the exact initialisation is $\mathcal{N}(0,\bI)$, and for stochastic interpolants the point-mass source $\delta_{\bx_0}$ is unchanged. The path-space construction of Appendix~\ref{appendix:proof_posterior} differs, propagating the tilt along the Markov SDE so that $\Phi_0^{\obs}$ is state-dependent and the reweighting is required.
\end{proof}

\begin{remark}[Degeneracy of the deterministic member for stochastic interpolants]
At $\gdiff_\tau=0$ the SDE reduces to the probability-flow ODE of the posterior path, transporting $p_0^{\obs}$ to $p_1^{\obs}$ deterministically. This needs a non-degenerate source: for flow matching $p_0=\mathcal{N}(0,\bI)$ and the map carries the posterior spread, whereas for stochastic interpolants $p_0=\delta_{\bx_0}$ is a point mass that the deterministic flow maps to a single point. The guided ODE is therefore available for flow matching but not for stochastic interpolants, whose randomness enters only through the diffusion term and is lost at $\gdiff_\tau=0$.
\end{remark}
